\documentclass[10pt]{article}

\usepackage[utf8]{inputenc}

\usepackage[T1]{fontenc}

\usepackage{amsmath}

\usepackage{amsfonts}

\usepackage{amssymb}

\usepackage[version=4]{mhchem}

\usepackage{stmaryrd}

\usepackage{graphicx}

\usepackage[export]{adjustbox}

\graphicspath{ {./images/} }



\begin{document}

$x^{(n-1)}(0)=\gamma_{n-1}$ автоматически приводится к алгебраич. уравнению и символически выражается формулой



\[

\begin{gathered}

x=\frac{\beta_{n-1} s^{n-1}+\ldots+\beta_{0}+f}{\alpha_{n} s^{n}+\ldots+\alpha_{0}}, \\

\beta_{v}=\alpha_{v+1} \gamma_{0}+\ldots+\alpha_{n} \gamma_{n-v-1},

\end{gathered}

\]



решение в обычном виде получается разложением на элементарные дроби от переменной $s$ с последующим обратным переходом по соответствующим таблицам $\boldsymbol{k}$ функциям.



Для применения О. и. к уравнениям с частными производными (а также к более общим псевдодифференциальным уравнениям) строятся дифференциальное и интегральное исчисления о п е р а т о р в ы х фу у кц и й, т. е. функций, значениями к-рых являются операторы: вводятся понятия непрерывности, производнӧ̈, сходимости ряда, интеграла и т. Д.



Пусть $f(\lambda, t)$ - нек-рая функция, определенная для $t \geqslant 0$ и $\lambda \in[a, b]$. П а р а м е т р и ч е с к а я о пер ат о р н а я фу у н ци я $f(\lambda)$ определяется формулой $f(\lambda)=\{f(\lambda, t)\}$; она ставит в соответствие рассматриваемым значениям $\lambda$ операторы частного вида - функции от $t$. Операторная функция наз. непрерывной при $\lambda \in[a, b]$, если она представима как произведение некрого оператора $q$ и такой параметрич. функции $f_{1}(\lambda)=$ $=\left\{f_{1}(\lambda, t)\right\}$, что $f_{1}(\lambda, t)$ непрерывна в обычном смысле. П р и м е р ы. 1) С помощью параметрич. функции $h(\lambda)=\{h(\lambda, t)\}$



\[

h(\lambda, t)=\left\{\begin{array}{cl}

0 & \text { для } 0 \leqslant t<\lambda, \\

t-\lambda & \text { для } 0 \leqslant \lambda \leqslant t,

\end{array}\right.

\]



определяется ф у н к и я $\mathrm{X}$ е в и с а їі д а



\[

H(\lambda)=s\{h(\lambda, t)\}

\]



значения г и п е р болич и с к о й п о к а з а т е льно Ї̈ ф у нкци и



\[

e^{-\lambda s} \equiv s H(\lambda)=s^{2}\{h(\lambda, t)\}

\]



наз. о п е р а то р а м и д в и г а, поскольку умножение данной функции на $e^{-\lambda s}$ вызывает смещение ее графика на длину $\lambda$ в положительном направлении оси $t$.



\begin{enumerate}

  \setcounter{enumi}{1}

  \item Решение уравнения теплопроводности

\end{enumerate}



\[

\frac{\partial x}{\partial t}=\alpha^{2} \frac{\partial^{2} x}{\partial \lambda^{2}}

\]



выражается через п а р а б о ли ч е с ку ю по к аз а т л л в у оф у нкци и (являюццуюся также параметрической операторноӥ функцией):



\[

e-\alpha \lambda \vee \bar{s}=\left\{\frac{\lambda}{2 \sqrt{\pi} t^{3}} \exp \left(-\frac{\lambda^{2}}{4 t}\right)\right\}

\]



\begin{enumerate}

  \setcounter{enumi}{2}

  \item Периодич. Функция $f(t)$ с периодом $2 \lambda_{0}$ имеет представление

\end{enumerate}



\[

\{f\}=\frac{\int_{0}^{2 \lambda_{0}} e^{-\lambda s_{f}(\lambda) d \lambda}}{1-e^{-2 \lambda_{0} s}} .

\]



\begin{enumerate}

  \setcounter{enumi}{3}

  \item Если $f(\lambda)$ принимает числовые значения в интервале $\left[\lambda_{1}, \lambda_{2}\right]$, то

\end{enumerate}



\[

\int_{\lambda_{1}}^{\lambda_{2}} e^{-\lambda s} f(\lambda) d \lambda=\left\{\begin{array}{cl}

f(\lambda), & \lambda_{1}<t<\lambda_{2}, \\

0, & 0 \leqslant t<\lambda_{1}, t>\lambda_{2},

\end{array}\right.

\]



т. е. умножение данной функции $\{f\}$ на $e^{-\lambda s}$ с последующим интегрированием вызывает у сечен и е ее графика. В частности,



\[

\int_{0}^{\infty} e^{-\lambda s} f(\lambda) d \lambda=\{f(t)\}

\]



такім образом, каждой функции $f(t)$, для к-рой рассмат- риваемый интеграл

аналитич. функция



\[

F(s)=\int_{0}^{\infty} e^{-s t} f(t) d t

\]



\begin{itemize}

  \item ее преобразование Лапласа. Благодаря этому обстоятельству довольно обширный класс операторов описывается функциями одного параметра $s$, более того, это формальное сходство уточняется математически установлением определенного изоморфизма.

\end{itemize}



Имеются различные обобщения О. и.; таково О. и. дифференциальных операторов, отличных от $s=\frac{d}{d t}$, напр. $b=\frac{d}{d t}\left(t \frac{d}{d t}\right)$, к-рое основывается на функциональных кольцах с надлежащим образом определенным произведением.



Лuт.: [1] Д и т к и н В. А., П р у Д н и ко в А. П., Справочник по операционному исчислению, М., 1965; [2] М и к у с и нМ. И. Войчеховский.



А-ОПЕРАЦИЯ, о п е р а п и я $A$, - теоретико-множественная операция, открытая П. С. Александровым [1] (см. также [2] с. 39 и [3]). Пусть $\left\{E_{n_{1} \ldots n_{k}}\right\}-$ система множеств, заиндексированных всеми конечными последовательностями натуральных чисел. Множество



\[

P=\bigcup_{n_{1} \ldots n_{k} \ldots k=1}^{\infty} \prod_{n_{1} \ldots n_{k}}

\]



где суммирование распространяется на все бесконечные последовательности натуральных чисел, наз. результатом $A-0 .$, примененной к системе $\left\{E_{n, \ldots n_{k}}\right\}$.



Применение $A$-О. к системе интервалов числовой прямой дает множества (названные $A$-множестважи в честь П. С. Александрова), к-рые могут не быть борелевскими (см. Дескриптивная теория множеств). А-О. сильнее операций счетного объединения и счетного пересечения и является идемпотентной. Относительно $A$-О. инвариантны $\boldsymbol{D}^{-}$. свойство (подмножеств произвольного топологич. пространства) и измеримость по Лебегу.



Лum.: [1] А л е к с а н д р о в П. С., "С. r. Acad. sci», 1916, t. 162, p. 323-25; [2] е г о же, Теория функций действительного переменного и теория топологических пространств, М.,

1978 ; [3] К о л м о г о о в А. Н., "Успехи матем. наук», 1966, т. 21, в. 4, с. 275-78; [4] С у с л и н М. Я., "С. r. Acad. sci.", 1917, t. $16 \frac{1}{4}$ p. 88 -91; [5] Л у з и н Н. Н., Собр. соч., т. 2, М.,



\includegraphics[max width=\textwidth, center]{2023_07_08_6028303e43c6bcab6ef2g-1}

$\boldsymbol{\delta} s$-OПЕРАЦИЯ - теоретико-множественная операция, результат применения к-рой к последовательности $\left(E_{n}\right)$ множкест может быть записан в виде



\[

\Phi\left(E_{n}\right)=\bigcup_{z \in N} \bigcap_{n \in z} E_{n},

\]



где $N$ - система множеств положительных целых чисел, наз. б а з й $\delta s=0$. См. Дескриптивная теория множеств.



Лит.: [1] К о л м о г о р о в А. Н., «Матем. сб.», 1928, т. 35, № 3-4, c. 415-22; [2] Х а у с д о $\mathrm{p}$ Ф Ф., Теория множеств, пер. с нем., М.- Ј., 1937; [3] А л е к с а н д р о в П. С., Теория функций действительного переменного и теория топологических

пространств, М., 1978, с. $35-40,53-58 ;[4] \mathrm{O}$ ч а ІО. С., «Успространств, М., 1978, c. $35-40,53-58 ;$ [4] О ч а н IO. С., "Ус-

пехи матем. наук», 1955 , т. 10 , в. 3 , с. $71-128 ;$ [ [5] К о л м о г ор о в А. Н., там же, 1966, т. 21 , в. 4, с. 275-78. А. Г. Елькин.



ОПЕРЕНИЕ П р о с т р а н с т в а - счетное семейство $P$ покрытий пространства $X$ множествами, открытыми в нек-ром объемлюцем пространстве $Y$, такое, что



\[

\cap\left\{\operatorname{St}_{\gamma}(x): \gamma \in P\right\} \subset X

\]



для каждої точки $x \in X$ (здесь $\operatorname{St}_{\gamma}(x)$ означает звезду точки $x$ относительно $\gamma$, т. е. объединение всех элементов из $\gamma$, содержащих точку $x$ ).



Понятие $\mathrm{O}$. лежит в основе определения т. н. pпространства (в смысле А. В. Архангельского). Про-





\end{document}